\section{Making the Torah Greek}

The Greek Pentateuch was probably produced in the middle or later third century \textsc{bc}. No colophon names its translators or gives a year. Differences among books and recurring agreements in vocabulary support a coordinated undertaking with more than one translator, but attempts to count hands precisely remain hypotheses. “Committee” is acceptable shorthand only if it does not import modern institutional procedures into the evidence.

\subsection{A translation language of its own}

The translators wrote intelligible Koine Greek, yet their work often follows Hebrew word order, coordination, repetition, and lexical patterning more closely than an independent Greek composition would. That texture is not simply bad Greek. It can mark reverence for the source, supply a stable correspondence for readers who knew both languages, and create a scriptural register whose unusual expressions became naturalized through repetition.

At the same time, the Greek Torah is not mechanically uniform. Translators transliterated some names, interpreted difficult expressions, selected Greek political and cultic terms, and sometimes made implicit relations explicit. Genesis can read differently from Leviticus; a rendering may be literal at the level of word order but interpretive in lexical choice. The usual line from “literal” to “free” is therefore too thin. Translation technique must ask what unit is being matched---morpheme, word, clause, discourse, or social institution---and what a Greek reader could understand.

The earliest project also created a lexicon for later translators. Greek equivalents for covenant, law, glory, righteousness, assembly, sacrifice, and named institutions became available for adoption, adaptation, or refusal. This is historical influence, not proof that every later translator worked in the same school or pursued one theology.

\subsection{The source behind the Greek}

The medieval Masoretic Text is indispensable for comparing Hebrew and Greek, but it is not the manuscript on the Alexandrian desk. When the Greek differs, at least five explanations must be tested:

\begin{enumerate}
\item the translator understood the same consonants differently;
\item the translator paraphrased, interpreted, harmonized, or made an error;
\item the translator's Hebrew exemplar contained a different reading;
\item the Hebrew and Greek represent different literary editions of a book; or
\item later Greek or Hebrew transmission created the difference.
\end{enumerate}

The existence of Hebrew textual plurality in the Judean Desert makes the third and fourth possibilities concrete. It does not make the Greek automatically earlier or better. Sometimes a Greek rendering transparently reflects interpretation or confusion; sometimes it preserves evidence for a Hebrew reading otherwise lost; sometimes no confident retroversion is possible. Textual criticism decides locally, not by awarding one tradition permanent victory.

\begin{fourstable}{Question}{Evidence consulted}{Responsible inference}{Overreach}
Did the translator have different consonants? & Greek syntax and vocabulary, surviving Hebrew variants, ancient versions, and the book's internal pattern. & A proposed source reading with stated confidence. & Retroverting every Greek difference into an unattested Hebrew manuscript.\\
Is the rendering interpretive? & Consistent equivalents, contemporary Greek usage, context, and analogous passages. & The Greek presents an ancient Jewish reading of its source. & Calling every intelligible difference a deliberate theological correction.\\
Which form is earliest? & Full transmission history and criteria appropriate to the variant. & One reading may better explain the others. & Declaring either MT or LXX universally “closest to the original.”\\
Is there one translation technique? & Book-level and unit-level profiles. & Related translators share some inherited vocabulary while differing in execution. & Treating a striking verse as the method of the whole corpus.\\
\end{fourstable}

\subsection{Use before preservation}

The Greek Torah was not only an artifact to be compared with Hebrew. It was read, copied, cited, interpreted, and made a basis for Jewish life in Greek. Aristeas's public ratification, Philo's exegesis, papyrus survival, and later vocabulary show authority and use. The exact place of Greek in synagogue reading during the originating century remains less securely documented than is sometimes claimed. Translation surely served access; its first precise ritual setting is not recoverable.

The achievement was nevertheless radical. A law bound to the history and worship of Israel could inhabit Greek without surrendering its source language. That possibility made every later stage of the Septuagint imaginable.
